\documentclass[11pt,a4paper]{article}
% ===== Packages =====
\usepackage[utf8]{inputenc}
\usepackage[T1]{fontenc}
\usepackage{times}      % Times text; math in Computer Modern
\renewcommand{\ttdefault}{cmtt}
\usepackage{microtype}
\usepackage{amsmath,amssymb,amsthm}
\usepackage{booktabs}
\usepackage{array}
\usepackage{graphicx}
\usepackage{xcolor}
\usepackage{tikz}
\usetikzlibrary{arrows.meta,positioning}
\usepackage{geometry}
\usepackage{xr-hyper}
\externaldocument{MDWF_MULTIGRID_NOTE}
\usepackage[colorlinks=true,linkcolor=blue!50!black,citecolor=green!35!black,urlcolor=blue!50!black]{hyperref}
\geometry{a4paper,left=2.5cm,right=2.5cm,top=2.4cm,bottom=2.6cm}
\setlength{\parindent}{1.2em}
\setlength{\parskip}{0pt}
\setlength{\emergencystretch}{2em}
\renewcommand{\baselinestretch}{1.05}

% ----- Algorithm environments -----
\newcounter{alg}
\newenvironment{algorithm}[1][]{%
  \par\medskip\noindent\begin{minipage}{\linewidth}\hrule\vspace{3pt}%
  \renewcommand{\caption}[1]{\refstepcounter{alg}\textbf{Algorithm~\thealg.} ##1\par\vspace{3pt}\hrule\vspace{2pt}}}%
  {\vspace{2pt}\hrule\end{minipage}\par\medskip}
\newenvironment{algorithmic}{%
  \begin{list}{\small\arabic{enumi}:}{\usecounter{enumi}\setlength{\topsep}{2pt}\setlength{\itemsep}{1pt}%
  \setlength{\leftmargin}{2.2em}\setlength{\labelwidth}{1.6em}}}{\end{list}}
\newcommand{\State}{\item}
\newcommand{\For}[1]{\item \textbf{for} #1 \textbf{do}}
\newcommand{\EndFor}{\item \textbf{end for}}
\newcommand{\Comment}[1]{\hfill{\small$\triangleright$ #1}}
\newcommand{\Return}{\textbf{return}\ }

% ----- Headings, abstract and captions -----
\makeatletter
\renewcommand\@seccntformat[1]{\csname the#1\endcsname.\hspace{0.6em}}
\renewcommand\thesection{\Roman{section}}
\renewcommand\thesubsection{\Alph{subsection}}
% A reference to a subsection names its section too ("Sec. IV.B", not "Sec. B").
\renewcommand\p@subsection{\thesection.}
\renewcommand\thesubsubsection{\arabic{subsubsection}}
\renewcommand\p@subsubsection{\thesection.\thesubsection.}
\renewcommand\section{\@startsection{section}{1}{\z@}%
  {3.0ex \@plus 1ex \@minus .2ex}{1.4ex \@plus .2ex}%
  {\centering\normalfont\normalsize\bfseries}}
\renewcommand\subsection{\@startsection{subsection}{2}{\z@}%
  {2.2ex \@plus 1ex \@minus .2ex}{1.1ex \@plus .2ex}%
  {\centering\normalfont\small\bfseries}}
\renewcommand\subsubsection{\@startsection{subsubsection}{3}{\z@}%
  {1.8ex \@plus 1ex \@minus .2ex}{0.9ex \@plus .2ex}%
  {\centering\normalfont\small\bfseries\itshape}}
\renewcommand\paragraph{\@startsection{paragraph}{4}{\z@}%
  {1.6ex \@plus 1ex \@minus .2ex}{-0.6em}%
  {\normalfont\normalsize\itshape}}
\renewenvironment{abstract}{%
  \begin{quote}\small\noindent}{\end{quote}\vspace{0.6ex}}
\renewcommand\maketitle{\begin{center}%
  {\large\bfseries\@title\par}\vspace{1.4ex}%
  {\normalsize\@author\par}\vspace{0.4ex}%
  {\small\itshape\@date\par}\end{center}\vspace{0.6ex}}
\renewcommand*\l@section{\@dottedtocline{1}{0pt}{2.4em}}
\long\def\@makecaption#1#2{\vskip 6pt{\small\textsc{#1.}~#2\par}}
\makeatother
\renewcommand{\figurename}{Fig.}
\renewcommand{\refname}{References}
\numberwithin{equation}{section}
\renewcommand{\theequation}{\arabic{section}.\arabic{equation}}

% ----- Statements -----
\theoremstyle{plain}
\newtheorem{proposition}{Proposition}
\theoremstyle{remark}
\newtheorem{remark}{Remark}
% Marker for statements that are derived here but not checked independently

% ===== Notation =====
\newcommand{\dg}{^{\dagger}}
\newcommand{\Dpv}{D_{\mathrm{PV}}}
\newcommand{\Cpv}{C_{\mathrm{PV}}}
\newcommand{\Spv}{S_{\mathrm{PV}}}
\newcommand{\Bpv}{B_{\mathrm{PV}}}
\newcommand{\Pp}{P_{+}}
\newcommand{\Pm}{P_{-}}
\newcommand{\Gfive}{\Gamma_5}
\newcommand{\AR}{\tilde A_{R}}
\newcommand{\AL}{\tilde A_{L}}
\newcommand{\Pro}{\mathcal P}
\newcommand{\Dc}{\mathcal D_{\mathrm c}}
\newcommand{\Dcpv}{\mathcal D_{\mathrm c}^{\mathrm{PV}}}
\newcommand{\Hc}{\mathcal H_{\mathrm c}}
\renewcommand{\Re}{\operatorname{Re}}
\newcommand{\Ce}{C_{e}}
\newcommand{\Co}{C_{o}}
\newcommand{\Cpve}{C_{\mathrm{PV},e}}
\newcommand{\Cpvo}{C_{\mathrm{PV},o}}

% Implementation and numerical-results sections; off = theory-only note.
\newif\ifresults
\resultstrue

\title{Reproducing the results of ``Multigrid for the even-odd preconditioned M\"obius
domain-wall operator with a factorized Pauli--Villars coarse correction''}
\author{Wei-Lun Chen}
\date{\today}

\begin{document}
\maketitle

\noindent This guide accompanies the paper of the title. It describes the implementation in
RutticeQCD and how to rerun the measurements of the paper and extend them. Section, table,
figure and equation numbers in it refer to the paper; the data and scripts are in
\texttt{docs/mg\_background} of the repository.

% How to reproduce the measurements and extend them to more configurations (input by
% MDWF_MULTIGRID_NOTE.tex with \resultstrue).

% ============================================================
\section{Reproducing and extending the measurements}
\label{sec:reproduce}
\begingroup\sloppy

All numbers of Sec.~\ref{sec:results} come from one configuration per volume; the
most useful extensions are more configurations, a larger volume and dynamical
ensembles. This section gives what is needed to run them.

\paragraph{Code.} RutticeQCD, \url{https://github.com/atsanroger/RutticeQCD}, branch
\texttt{cuda-bridge} (the commit of each result is given with it; later commits give
the same iterations, up to the two-iteration differences between builds noted in
Sec.~\ref{sec:massscan}). It needs Rust (stable) and, for the GPU, CUDA~12:
{\small
\begin{verbatim}
  export CUDA_HOME=/usr/local/cuda
  cargo build --release --features cuda-bridge
  ./target/release/ruttice mg_24.yaml
\end{verbatim}}
The GPU backend runs on one GPU. The parameter files are in the style of Bridge++
and are documented in \texttt{docs/USAGE.md}; gauge configurations are read in the
Bridge++ text and binary formats, NERSC, ILDG and SciDAC.

\paragraph{The mass scan (Table~\ref{tab:massscan}, Fig.~\ref{fig:massscan}).} The
parameter files are \texttt{docs/mg\_background/runs/mg\_24.yaml} (the multigrid) and
\texttt{runs/cgne\_24.yaml} (the reference, double or mixed precision); the comments
in them say how to switch to setup~B and to a point source. Each run prints a line
{\small
\begin{verbatim}
  seed 91: GCR 26 iterations on the device; 2875 coarse applications;
      worst |Dx-b|^2 2e-18; 24.9 s; ... device 19.91 GB
\end{verbatim}}
and, for the multigrid, \texttt{setup on the device: 393 s}. These, per mass and
setup, are what Table~\ref{tab:massscan} lists. On an A30 one configuration takes about
$1.5$ hours for the multigrid (seven masses, both setups; the setup dominates) and
about $40$ minutes for the two CGNE scans.

\paragraph{Memory and volume.} On $24^3\times48$ the multigrid needs $19.9$\,GB of
device memory (CGNE $6.9$\,GB in double, $10.1$\,GB mixed), roughly proportional to the
volume: $32^3\times64$ needs about $63$\,GB (an $80$\,GB GPU), $48^3\times96$ more
than one GPU has. The aggregates, $4^4$, must divide the lattice, and each level
below the coarse lattice must divide the extent of the level above
(\texttt{MultiGrid\_Level2}, \texttt{\_Level3}; on $32^3\times64$, for instance,
super-blocks $2$ and $2$).

\paragraph{The spectral quantities (Secs.~\ref{sec:fov}--\ref{sec:pseudo}).} The
example \texttt{preconditioned\_fov\_device} computes, with exact coarse solves, the
field of values, the parts of the cycle, the eigenvalues with their condition
numbers, and $\sigma_{\min}(z-\tilde A\mathcal M)$ on contours, e.g.
{\small
\begin{verbatim}
  cargo run --release --features cuda-bridge \
      --example preconditioned_fov_device -- \
      16 16 binary conf.bin 1.0 0.001 1 --ratio 10 \
      --arnoldi --far 1,0 --pseudo --contour 1,0,0.6,16
\end{verbatim}}
(arguments: $L$, $T$, format, file, $M_5$, $m$, stout $0/1$). It fits $16^4$ in $24$\,GB
and takes about an hour there; $24^3\times48$ does not fit. The CPU examples
\texttt{preconditioned\_fov} and \texttt{coarse\_hermiticity} need no GPU.

\paragraph{What would help most.}
\begin{enumerate}
\item The mass scan of Table~\ref{tab:massscan} on three to five further
$24^3\times48$ configurations, for the spread of the iterations and of the
speed-up.
\item The same on a larger volume ($32^3\times64$ on an $80$\,GB GPU), for the volume
dependence of the multigrid iterations and of the setup.
\item The mass scan on a dynamical ($2+1$-flavor) ensemble, e.g.\ that of
Ref.~\cite{KCM2024}, where the spectra of Sec.~\ref{sec:spectra} were taken.
\item The spectral run above on further $16^4$ configurations, for how the outliers
of setup~S and the pseudospectra vary between configurations.
\end{enumerate}
For each run, the result lines above and the commit of the build are enough; the
logs contain everything else.
\endgroup


\begin{thebibliography}{9}
\bibitem{KCM2024} I.~Kanamori, W.-L.~Chen and H.~Matsufuru, \emph{Spectrum of preconditioned Moebius domain-wall operators}, PoS \textbf{LATTICE2024} (2025) 414.
\bibitem{Stewart} G.~W.~Stewart, \emph{A Krylov--Schur algorithm for large eigenproblems},
SIAM J.\ Matrix Anal.\ Appl.\ \textbf{23} (2001) 601.
\end{thebibliography}
\end{document}
